\section{Self-Consistency: del muestreo a la votación}

\subsection{Motivación: la naturaleza probabilística de los LLM}

Denny Zhou parte de primeros principios para derivar la validez de Self-Consistency:

\begin{itemize}
    \item Al decodificar, el LLM optimiza: $\arg\max P(\text{reasoning path}, \text{answer} \mid \text{problem})$
    \item Pero lo que en realidad necesitamos es: $\arg\max P(\text{answer} \mid \text{problem})$
    \item La relación entre ambos: $P(\text{answer} \mid \text{problem}) = \sum_{\text{path}} P(\text{answer}, \text{path} \mid \text{problem})$
    \item Es decir, hay que sumar (marginalizar) sobre todas las rutas de razonamiento posibles, en lugar de tomar solo la única ruta óptima.
\end{itemize}

\subsection{Método y resultados}

\begin{importantbox}{Método de Self-Consistency}
\begin{enumerate}
    \item Muestrear el mismo problema varias veces (con temperatura distinta de cero) para obtener varias rutas de razonamiento y sus respuestas correspondientes.
    \item Tomar la \textbf{respuesta con mayor frecuencia de aparición} como predicción final (votación mayoritaria / majority voting).
    \item Nota: lo que se vota es la \textbf{respuesta final}, no la ruta de razonamiento; la ruta de razonamiento es una variable latente (latent variable).
\end{enumerate}
\end{importantbox}

Self-Consistency trajo una enorme mejora de desempeño y en su momento superó ampliamente el estado del arte en todos los puntos de referencia. Más importante aún, cuando la consistencia supera el 80\%, la exactitud se acerca al 100\%, lo cual ofrece un indicador confiable de confianza.

\begin{knowledgebox}{Extensión con Universal Self-Consistency}
Para respuestas de forma libre (free-form), el propio LLM puede juzgar qué respuestas son semánticamente equivalentes y así encontrar el grupo de respuestas más común. Por ejemplo, para una pregunta como «qué países tienen un consumo per cápita de café menor que el de México», distintas respuestas pueden tener una redacción diferente pero un contenido consistente.
\end{knowledgebox}

\subsection{Resumen de la sección}

Self-Consistency eleva el razonamiento del LLM de una única generación al nivel de la inferencia estadística. Su idea central proviene de la inferencia marginal (marginal inference) de los modelos gráficos probabilísticos: es simple pero sumamente efectiva.

